\section*{Capítulo \ref{ch:b3:structural-biology} --- Biologia estrutural das proteínas}

\begin{solution}{exo:b3:structural-biology:1}
Primária: a sequência dos $141$ resíduos de uma cadeia $\alpha$.
Secundária: suas oito $\alpha$-hélices. Terciária: o \omterm{def:b3:structural-biology:levels}{enovelamento} de
globina de uma cadeia envolvendo seu heme. Quaternária: o tetrâmero
$\alpha_{2}\beta_{2}$ e o modo como seus dímeros giram um contra o
outro.
\end{solution}

\begin{solution}{exo:b3:structural-biology:2}
$36\times \qty{0.15}{nm} = \qty{5.4}{nm}$; $36/3.6 = 10$ voltas;
$36 - 4 = 32$ pontes de hidrogênio.
\end{solution}

\begin{solution}{exo:b3:structural-biology:3}
Que a estrutura nativa, inclusive o pareamento correto de quatro pontes
dissulfeto entre $105$ possibilidades, é recuperada espontaneamente a
partir da cadeia desenovelada: a sequência contém toda a informação
necessária, e o estado nativo é o mínimo termodinâmico. Fazê-lo num
tubo de ensaio com proteína pura excluiu qualquer molde, enzima ou
maquinaria celular como fonte da informação.
\end{solution}

\begin{solution}{exo:b3:structural-biology:4}
RMN: proteínas pequenas, abaixo de uns \qty{40}{kDa}, em solução.
Cristalografia: qualquer tamanho que cristalize, de peptídeos a
ribossomos. \omterm{def:b3:structural-biology:methods}{Criomicroscopia eletrônica}: complexos grandes, acima de uns
\qty{50}{kDa}, sem cristais. A RMN informa diretamente sobre o
movimento (e a crio-ME consegue separar as moléculas em classes
conformacionais); uma estrutura cristalográfica é uma média estática.
\end{solution}

\begin{solution}{exo:b3:structural-biology:5}
$2^{150} \approx 1.4\times 10^{45}$ conformações, $1.4\times 10^{33}$
s $\approx 5\times 10^{25}$ anos. $3^{20} \approx 3.5\times 10^{9}$,
$3.5$ ms: o peptídeo poderia fazer uma busca exaustiva dentro de um
segundo; a proteína não conseguiria na idade do universo.
\end{solution}

\begin{solution}{exo:b3:structural-biology:6}
$Y = [\alpha(1+\alpha) + Lc\alpha(1+c\alpha)]/[(1+\alpha)^{2} +
L(1+c\alpha)^{2}]$. $\alpha = 1$: $3.01/106 = 0.028$ (sítio único
$0.50$). $\alpha = 10$: $121/242 = 0.50$ ($0.91$). $\alpha = 100$:
$\num{10300}/\num{10601} = 0.97$ ($0.99$). A curva MWC fica bem atrás
do sítio único em ligante baixo e depois sobe abruptamente --- de
$0.03$ a $0.97$ ao longo de duas décadas em que o sítio único vai de
$0.5$ a $0.99$: essa inclinação é a \omterm{def:b3:structural-biology:allostery}{cooperatividade}.
\end{solution}

\begin{solution}{exo:b3:structural-biology:7}
O bisfosfoglicerato liga apenas o \omterm{def:b3:structural-biology:allostery}{estado T} (na cavidade central), e por
isso estabiliza o T: no modelo ele eleva $L$, a razão $T/R$ da proteína
sem ligante, e assim desloca a curva para a direita --- afinidade menor,
mais oxigênio liberado às pressões teciduais. As hemácias armazenadas
perdem o composto, $L$ cai, a curva se desloca para a esquerda, e o
sangue transfundido retém seu oxigênio em vez de cedê-lo aos tecidos,
até que as células o ressintetizem ao longo de um dia.
\end{solution}

\begin{solution}{exo:b3:structural-biology:8}
$d = \lambda/(2\sin\theta) = 0.1/(2\times 0.259) = \qty{0.19}{nm}$: o
esqueleto e cada cadeia lateral são posicionados, com os átomos
individuais resolvidos. Para \qty{0.12}{nm}: $\sin\theta = 0.1/0.24 = 0.417$, $\theta =
\qty{25}{\degree}$.
\end{solution}

\begin{solution}{exo:b3:structural-biology:9}
Uma lisina carregada enterrada entre cadeias laterais apolares custa
dezenas de quilojoules por mol (sua carga sem solvatação, perdido o
enterramento hidrofóbico da leucina), mais do que todo o $\Delta G$ do
\omterm{def:b3:structural-biology:levels}{enovelamento}: a proteína se desenovela ou o cerne se rearranja. Na
superfície a lisina fica hidratada como ficaria no estado desenovelado
e nada se perde. $\qty{-30}{kJ/mol} \approx -12RT$ significa
$K = e^{12} \approx 1.6\times
10^{5}$: uma molécula em \num{160000}
está desenovelada a cada momento, e toda a margem equivale a duas ou
três pontes de hidrogênio entre centenas --- o valor de uma mutação.
\end{solution}

\begin{solution}{exo:b3:structural-biology:10}
As espécies $R$ somam $[R_{0}](1+\alpha)^{n}$ e as espécies $T$,
$[R_{0}]L(1+c\alpha)^{n}$, de modo que
$\bar R = (1+\alpha)^{n}/[(1+\alpha)^{n}
+ L(1+c\alpha)^{n}]$. Para
$c \to 0$ o termo $T$ é $L$, e $\bar R =
1/2$ quando
$(1+\alpha)^{n} = L$, isto é, $\alpha = L^{1/n} - 1$. Para a
hemoglobina $(3\times 10^{5})^{1/4} = 23.4$, logo $\alpha \approx 22$:
\qty{22}{mmHg}, perto do $P_{50}$ de \qty{26}{mmHg} --- a troca
conformacional e a meia-saturação quase coincidem.
\end{solution}

\begin{solution}{exo:b3:structural-biology:11}
Ela foi resistida porque todo agente infeccioso conhecido se replicava
por ácido nucleico, e uma proteína não consegue copiar sua sequência.
Encontrar um ácido nucleico necessário, ou mostrar que proteína
purificada não transmite, a refutaria. Apoio: a infecciosidade
sobrevive a nucleases e à radiação ultravioleta e ionizante em doses
que destroem qualquer genoma, mas é destruída por desnaturantes de
proteína e por proteases; a PrP recombinante enovelada in vitro em
fibrilas transmite a doença a animais, com as mesmas propriedades de
cepa ao longo das passagens. É preciso uma semente porque a conversão
de um monômero normal é proibitivamente lenta sozinha --- o núcleo da
estrutura cross-$\beta$ precisa formar-se primeiro, e esse é o evento
limitante ---, ao passo que a ponta de uma fibrila já formada converte
monômeros depressa: a forma é copiada por crescimento sobre molde.
\end{solution}

\begin{solution}{exo:b3:structural-biology:12}
Muitas sequências se enovelam na mesma estrutura: o que um \omterm{def:b3:structural-biology:levels}{enovelamento}
exige é um padrão de posições hidrofóbicas e polares, algumas glicinas
e prolinas nas curvas e os resíduos que fazem o trabalho; todo o resto
na superfície pode mudar sem penalidade, de modo que as mutações se
acumulam ali enquanto a seleção preserva o \omterm{def:b3:structural-biology:levels}{enovelamento} e a função. Os
resíduos conservados são, portanto, o cerne (mais como padrão do que
como identidades), os resíduos catalíticos e de ligação, e as glicinas,
prolinas e cisteínas estruturais. Os métodos de \omterm{def:b3:bioinformatics:msa}{perfil} ponderam
exatamente essas colunas e por isso reconhecem membros da família a
\qty{15}{\%} de identidade; o projeto de proteínas percorre a lógica ao
contrário, escolhendo um padrão hidrofóbico para um enovelamento-alvo e
deixando a superfície ser qualquer coisa.
\end{solution}

\begin{solution}{pb:b3:structural-biology:1}
\textbf{1.} $0.75\times 150 \approx 112$ resíduos helicoidais; $112\times
\qty{0.15}{nm} = \qty{17}{nm}$; $112/3.6 \approx 31$ voltas.
\textbf{2.} $112 - 8\times 4 = 80$ pontes de hidrogênio.
\textbf{3.} Massa $150\times 110 = \qty{16500}{Da} = 2.7\times 10^{-20}$
g; volume $2.7\times 10^{-20}\times 0.73 = 2.0\times 10^{-20}$ cm$^{3}$
$= \qty{20}{nm^3}$; raio $(3V/4\pi)^{1/3} = \qty{1.7}{nm}$.
\textbf{4.} Estendida, $150\times 0.36 = \qty{54}{nm}$ contra um
diâmetro de \qty{3.4}{nm}: um fator de $16$.
\textbf{5.} Os resíduos helicoidais são \qty{75}{\%} do comprimento
estendido, $112\times 0.36 = \qty{40}{nm}$; a hélice os comprime a
\qty{17}{nm}, um fator $0.36/0.15 = 2.4$.
\textbf{6.} $40\times 4 = \qty{160}{kJ/mol}$ ganhos, $150\times 0.9 =
\qty{135}{kJ/mol}$ perdidos: no total,
$\Delta G \approx \qty{-25}{kJ/mol}$, dentro da faixa marginal.
\textbf{7.} $3^{150} \approx 4\times 10^{71}$; $4\times 10^{59}$ s
$\approx 10^{52}$ anos.
\textbf{8.} $10^{-2}\times 10^{12} = 10^{10}$ conformações, uma fração
$3\times 10^{-62}$ da contagem.
\textbf{9.} Hélices em \qty{1}{\micro s} (em paralelo); $8! = \num{40320}$
empacotamentos a $10^{6}$ por segundo levam \qty{40}{ms}: no total,
cerca de \qty{40}{ms}, a ordem de grandeza observada. Sim: resolver as
partes de modo independente e depois o todo é um funil --- o espaço de
busca é fatorado, e não esgotado.
\textbf{10.} Ciclos esperados $1/0.3 = 3.3$, cerca de \qty{33}{s}; após
seis ciclos, $0.7^{6} = 0.12$ continua desenovelado.
\textbf{11.} $\Delta G = -RT\ln K$ com $K = 10^{4} \to 10^{2}$: de
$-23.7$ a \qty{-11.9}{kJ/mol}, uma perda de \qty{12}{kJ/mol} de
estabilidade ($RT = \qty{2.58}{kJ/mol}$, $\ln 100 = 4.6$).
\textbf{12.} A agregação começa por um núcleo, cuja taxa de formação
sobe como uma potência alta da concentração da espécie desenovelada
propensa à agregação; cem vezes mais dessa espécie eleva em ordens de
grandeza a chance de um núcleo se formar dentro de uma vida, e, uma vez
existindo o núcleo, o crescimento é rápido e se molda a si mesmo.
\textbf{13.} $\alpha = 100$: $Y = 1.054\times 10^{8}/1.089\times 10^{8} =
0.97$. $\alpha = 20$: numerador $185\,220 + 103\,680 = 288\,900$,
denominador $194\,481 + 622\,080 = 816\,561$: $Y = 0.35$.
\textbf{14.} $\bar R = 1.041\times 10^{8}/1.089\times 10^{8} = 0.96$ no
pulmão; $194\,481/816\,561 = 0.24$ no músculo.
\textbf{15.} $0.97 - 0.35 = 0.62$ da capacidade descarregada (dois
terços da carga). Sítio único: $100/126 - 20/46 = 0.79 - 0.43 = 0.36$.
\textbf{16.} Capacidade $150\times 1.34 = \qty{200}{mL/L}$; ao músculo,
$0.62\times 200 \approx \qty{125}{mL/L}$. Em repouso,
$(0.97 - 0.78)\times
200 = \qty{38}{mL/L}$; $250/38 \approx \qty{6.5}{L/min}$ de fluxo sanguíneo --- a ordem de um débito cardíaco
de repouso.
\textbf{17.} $L = 3\times 10^{6}$, $\alpha = 40$: numerador $2.76\times
10^{6} + 3.29\times 10^{6} = 6.05\times 10^{6}$, denominador $2.83\times
10^{6} + 11.5\times 10^{6} = 14.4\times 10^{6}$: $Y = 0.42$ em vez de
$0.78$ --- bem mais oxigênio descarregado em repouso, a adaptação do
sangue à altitude e do músculo ao trabalho.
\textbf{18.} Fetal, $L = 3\times 10^{4}$, $\alpha = 30$: numerador
$893\,730 + 19\,770 = 913\,500$, denominador $923\,520 + 85\,680 =
1\,009\,200$: $Y = 0.91$. Materna à mesma pressão: numerador
$893\,730 + 197\,730 = 1\,091\,460$, denominador $923\,520 + 856\,830
= 1\,780\,350$: $Y = 0.61$. À mesma pressão o sangue fetal retém mais
oxigênio, e por isso o oxigênio flui da mãe para o feto através da
placenta.
\textbf{19.} $(10^{5}\,\text{nm}/6\,\text{nm})^{3} = 4.6\times 10^{12}$
células; $1.9\times 10^{13}$ moléculas.
\textbf{20.} $d = 0.1/(2\sin 14.5^{\circ}) = \qty{0.20}{nm}$: sim, cada
cadeia lateral e cada átomo ficam posicionados.
\textbf{21.} $\sin\theta = 0.1/0.3$, $\theta = \qty{19.5}{\degree}$. Num
cristal desordenado as moléculas não estão em registro exato, as ondas
espalhadas a ângulos altos --- que codificam os detalhes finos --- deixam
de somar-se em fase, e os pontos se apagam em ângulos altos; só
sobrevivem os pontos de baixa \omterm{def:b3:structural-biology:methods}{resolução}.
\textbf{22.} Reduzir $d$ à metade exige o dobro do sinal e, portanto,
quatro vezes as partículas: cerca de \num{40000} (mais, na prática, já
que outros erros entram).
\textbf{23.} Proteínas de membrana, que raramente cristalizam a partir
de soluções de detergente e podem ser fotografadas num nanodisco
lipídico; complexos grandes e flexíveis --- ribossomos, spliceossomos,
canais iônicos com seus reguladores ---, cuja heterogeneidade impede a
cristalização, mas cujas partículas podem ser separadas em classes
conformacionais.
\textbf{24.} O heme e o oxigênio ligados e sua geometria exata, as
águas e os íons ordenados, a segunda conformação (T contra R) e a prova
de que o \omterm{def:b3:structural-biology:levels}{enovelamento} está certo --- uma predição é um único modelo sem
ligante e sem conferência experimental.
\textbf{25.} Descarga de $0.62$ da capacidade (sítio único $0.36$); o
tempo de Levinthal, cerca de $10^{52}$ anos; um mapa a \qty{0.20}{nm}.
\end{solution}
